\scanpage{60}
15) \uline{Cohomologie absolue}

En plus des foncteurs cohomologiques de motifs $\mathbb R^{\unsure{*|+}}f_*$ [$f$ morphisme de type fini] il y a lieu de considérer également des foncteurs ``absolus''
\[
H^*(X,M)=(H^i(X,M))_i
\]
$M$ motif variable sur $X$, qui doivent provenir d'un foncteur de catégories triangulées
\[
\mathbb R\Gamma_X:\ D^b(\deleted{\unsure{\mathcal M}}\mathcal M^{\deleted{\unsure{+}}}(X))\longrightarrow D^b(\mathrm{Ab})
\]
On veut que les $H^*(X,M)$ et $\mathbb R\Gamma(\unsure{M.|M_.|M_{\bullet}})$ aient les propriétés suivantes

$1^\circ)$ Commutation aux $\varprojlim$ essentiellement affines de préschémas . \unsure{qu.|q.} cpt \unsure{qu.|q.} sép. Cela nous ramène pratiquement aux cas $X$ de type fini sur $\mathbb Z$..

$2^\circ)$ Transitivité : \unsure{Si|si} $f:X\mathbin{-}Y$ de
% archive sheet number: 59

